\documentclass[a4paper]{article}

% --- Packages ---
\usepackage{amsmath, amssymb} % Equations and math symbols
\usepackage{graphicx}         % Figures
\usepackage{listings}         % Code blocks
\usepackage{cite}             % BibTeX citations
\usepackage{hyperref}         % Clickable references
\usepackage{labeltext}

% --- Code Block Configuration ---
\lstset{
    basicstyle=\ttfamily\footnotesize,
    breaklines=true,
    numbers=left,
    numberstyle=\tiny,
}

\title{Reassessment of the microTEDAclus Clustering Algorithm and Proposal for an Improved Model}
\author{Eduardo Henrique Basilio de Carvalho \\ Advisor: Prof. Janier Arias Garcia}
\date{\today}

\begin{document}

\maketitle

\begin{abstract}
    The abstract is deferred to after we finish the main content of the document. This allows us to write a more informed and concise summary of the work after all sections have been completed.
\end{abstract}

\section{Introduction}
The introduction is also deferred until we have a clearer picture of the entire document. We will revisit this section after we have completed the main body of the report to ensure it effectively sets the stage for our research and findings.

\section{Background and Related Work}

The TEDA framework is presented as an alternative to probability theory for handling data in non-deterministic scenarios \cite{angelov2014outside}. The original microTEDAclus model applies its metrics to decide if a sample is an outlier to a microcluster and if microclusters are suitable candidates to form a macrocluster \cite{maia2020evolvclus}. Finally, when evaluating the previous model's performance in the last-mile problem, \cite{galindo2021lastmile} refactor the decay logic and variance computation.

\subsection{Typicality and Eccentricity-Based Data Analytics}

The paper \textit{Outside the Box: An Alternative Data Analytics Frame-Work}, by Plamen Angelov, introduces the mathematical formulation for the metrics \textbf{Eccentricity} and \textbf{Typicality} \cite{angelov2014outside}.

Let

\begin{equation}
    C = \{ \sum_{i=1}^{k} \pi^k_i > 0,\quad j > 1,\quad k > 2 \}
\end{equation}

be a set of conditions for a sample $j$ to be considered for the computation of eccentricity and typicality in a moment $k$.

Eccentricity, a measure of how much a sample deviates from a cluster, is defined as:

\begin{equation}
    \xi^k_j = \frac{2\pi^k_j}{\sum_{i=1}^{k} \pi^k_i}, \quad C
\end{equation}

Where $\pi^k_j$ is the sum of distances from the sample $j$ to all samples in the moment $k$.

Typicality, a measure of how well a sample fits within a cluster, is defined as:

\begin{equation}
    \tau^k_j = 1 - \xi^k_j, \quad C
\end{equation}

Both are, by definition, bounded: $0 < \xi^k_j < 1$ and $0 < \tau^k_j < 1$. The normalized eccentricity and typicality are presented as:

\begin{equation}
    \zeta^k_j = \frac{\xi^k_j}{2}, \quad \sum_{i=1}^{k} \zeta^k_i = 1\quad 0 < \zeta^k_j < \frac{1}{2}\quad C
\end{equation}

\begin{equation}
    t^k_j = \frac{\tau^k_j}{k - 2}, \quad \sum_{i=1}^{k} t^k_i = 1\quad 0 < t^k_j < \frac{1}{k - 2}\quad C
\end{equation}

The paper states that samples for which the normalized eccentricity $\zeta^k_j > 1/k$ are considered anomalous, and those for which the normalized typicality $t^k_j > 1/k$ are considered typical.

Furthermore, it provides recursive, or incremental, formulas for updating $\sum_{i=1}^{k} \pi^k_i$ in the squared Euclidean distance, as follows:

\begin{gather}
    \sum_{i=1}^{k} \pi^k_i = \sum_{i=1}^{k-1} \pi^{k-1}_i + 2 \pi^k_k \\
    \pi^k_j = k (\lVert x_j - \mu^k \rVert^2 + X^k - \lVert \mu^k \rVert^2) \\
    \mu^k = \frac{k - 1}{k} \mu^{k-1} + \frac{1}{k} x_k \\
    X^k = \frac{k - 1}{k} X^{k-1} + \frac{1}{k} \lVert x_k \rVert^2
\end{gather}

Finally, it is made clear \labeltext[Angelov's statement about the use of other proximity functions]{those metrics suit various proximity functions, among which are the cosine similarity and the Mahalanobis distance}{info:mahalanobis_as_teda_metric} \cite{mahalanobis2018generalized}.

\subsection{Original microTEDAclus}



\subsubsection{Critique}

\subsection{microTEDAclus Updates in the \textit{Mineração de Dados Aplicada a Problemas de Last-Mile na Logística} PHD Thesis}

\subsubsection{Critique}

\section{Reassessment}

\subsection{Data}

\subsection{Preprocessing}

\subsection{Processing}

\section{Proposal for an Improved Model}

% --- Bibliography ---
\bibliographystyle{plain}
\bibliography{bibliography}

\end{document}
